\chapter{I/O and Testing}

\textbf{Readings:} CP:AMA 2.5, 3.1--3.2, 22.1--22.3 (selected).

\section{Input and Output}

\begin{definition}
\emph{Input and output} (I/O) is how a program interacts with the outside world: receiving data
(input) and displaying or sending data (output).
\end{definition}

A program may interact with humans (keyboard, screen) or with non-human entities: files,
other programs, or other computers over a network (e.g.\ a web server). In CS 136 we only use
\emph{standard input} (stdin, usually the keyboard) and \emph{standard output} (stdout, usually
the screen). Either can be \emph{redirected} to come from or go to a file, which is how Seashell
and Marmoset test our programs.

\section{Output with \code{printf}}

We have already used \code{printf} with the placeholders \code{\%d}, \code{\%c}, \code{\%f},
\code{\%zd} and \code{\%p} (and \code{\%s} for strings in Unit~08). Some details:
\begin{itemize}
  \item The number of placeholders must match the number (and types) of extra arguments.
    A mismatch is undefined behaviour, but the compiler usually warns.
  \item \code{printf} produces output as a side effect. (It also returns the number of
    characters printed, which we rarely use.)
  \item Output is often \emph{buffered}: it may not appear until a newline is printed or the
    program ends. If a program crashes, the last few lines of output can be lost, which can be
    misleading when debugging.
\end{itemize}

\section{Debugging Output and Tracing}

\emph{Tracing} means printing intermediate values to see what a program is doing. To avoid
having to delete tracing code before submitting, wrap it in a flag:
\begin{lstlisting}[style=c]
const bool DEBUG = true;   // set to false before submitting

int sum_first(int n) {
  int sum = 0;
  for (int i = 1; i <= n; ++i) {
    sum += i;
    if (DEBUG) {
      printf("TRACE: i = %d, sum = %d\n", i, sum);
    }
  }
  return sum;
}
\end{lstlisting}
\begin{aside}
A preprocessor version (\code{\#ifdef DEBUG ... \#endif}) removes the code completely at compile
time. Printing traces to \code{stderr} with \code{fprintf(stderr, ...)} keeps them separate from
the program's real output.
\end{aside}

\section{Input with \code{scanf}}

\code{scanf} reads formatted input. It uses placeholders like \code{printf}, but each extra
argument must be a \emph{pointer} to the variable that will receive the value.
\begin{lstlisting}[style=c]
int n = 0;
int count = scanf("%d", &n);   // note the &
\end{lstlisting}

\begin{important}[The return value of \code{scanf}]
\code{scanf} returns the number of values successfully read and stored:
\begin{itemize}
  \item \code{1} (for one placeholder): success;
  \item \code{0}: the input was not in the expected format (e.g.\ a letter when reading an int);
    the offending characters are left unread;
  \item \code{EOF} (a negative constant): the end of the input was reached before anything
    could be read.
\end{itemize}
Always check the return value. Never assume \code{scanf} succeeded.
\end{important}

\emph{End of file} (EOF) is a special signal that there is no more input. When input comes from
a file it happens at the end of the file; when typing in Seashell you can send it with
\textbf{Ctrl+D}.

\subsection{Whitespace}

\begin{itemize}
  \item \code{\%d} (and \code{\%f}, \code{\%s}) skip leading whitespace (spaces, tabs, newlines)
    automatically.
  \item \code{\%c} does \emph{not}: it reads the very next character, which may be a space or
    newline. To skip whitespace before reading a character use \code{" \%c"} (with a space).
\end{itemize}

\subsection{Reading all integers until EOF}

\begin{lstlisting}[style=c]
// read_sum() reads integers from input until EOF (or invalid input)
//   and produces their sum
// effects: reads input
int read_sum(void) {
  int sum = 0;
  int n = 0;
  while (scanf("%d", &n) == 1) {
    sum += n;
  }
  return sum;
}
\end{lstlisting}

\subsection{Reading commands}

Many assignments read a sequence of commands. A typical structure:
\begin{lstlisting}[style=c]
int main(void) {
  char cmd = 0;
  int arg = 0;
  while (scanf(" %c", &cmd) == 1) {
    if (cmd == 'a') {            // "a 5" adds 5
      if (scanf("%d", &arg) != 1) {
        printf("invalid input\n");
        break;
      }
      // ... add arg ...
    } else if (cmd == 'p') {     // "p" prints
      // ... print ...
    } else if (cmd == 'q') {     // "q" quits
      break;
    } else {
      printf("unknown command %c\n", cmd);
    }
  }
}
\end{lstlisting}

\subsection{Reading characters with \code{getchar}}

\code{getchar()} (in \code{<stdio.h>}) reads one character and returns it as an \code{int}, or
\code{EOF}. The result must be stored in an \code{int}, not a \code{char}, so that \code{EOF} can
be distinguished from every valid character.
\begin{lstlisting}[style=c]
// count_lines() counts the lines of input
int count_lines(void) {
  int count = 0;
  int c = 0;
  while ((c = getchar()) != EOF) {
    if (c == '\n') ++count;
  }
  return count;
}
\end{lstlisting}

\subsection{Invalid input}

When \code{scanf("\%d", \&n)} returns 0, the bad characters are still waiting in the input, so
calling it again in a loop would loop forever. Either stop, or skip the bad input (e.g.\ read
characters with \code{getchar} until a newline). Unless an assignment says otherwise, you may
assume input is valid, but your loops should still terminate on EOF.

\section{Testing}

\subsection{Why test?}

Testing cannot prove a program correct, but good tests find most bugs. In CS 136, public tests
on Marmoset only check basic functionality; private tests are much more thorough, so your own
testing determines your mark.

\subsection{Kinds of testing}

\begin{itemize}
  \item \textbf{Black-box testing:} tests are written from the specification (documentation)
    alone, without looking at the code. Good for checking the interface is correct.
  \item \textbf{White-box testing:} tests are written while looking at the code, to make sure
    every branch, loop and special case is exercised.
  \item \textbf{Assertion-based testing:} a test client calls functions and \code{assert}s the
    results (no I/O required). Good for modules.
  \item \textbf{I/O-driven testing:} the program reads input and produces output; tests compare
    the actual output with the expected output.
  \item \textbf{Regression testing:} re-running all old tests after every change to make sure
    nothing that used to work is now broken.
\end{itemize}

\subsection{What to test}

\begin{itemize}
  \item \textbf{Boundary / edge cases:} 0, 1, negative numbers, \code{INT\_MAX}, empty input,
    an empty collection, a collection with exactly one item, the first and last positions.
  \item \textbf{Each branch} of every conditional, and loops that run 0, 1 and many times.
  \item \textbf{Typical cases} (to check general correctness).
  \item \textbf{Large cases} (to check efficiency and stack depth).
  \item \textbf{Mutation:} check the state after each side effect, not only the returned value.
\end{itemize}

\subsection{Testing in Seashell}

For a program \code{prog.c} that reads input, create a pair of files in the \code{tests}
folder:
\begin{itemize}
  \item \code{name.in}: the input to feed to the program;
  \item \code{name.expect}: the exact expected output.
\end{itemize}
Running tests feeds each \code{.in} file to the program and reports whether the output matched
the \code{.expect} file exactly (whitespace matters, including the final newline). This is the
same mechanism Marmoset uses. Outside Seashell you would do the same with redirection:
\begin{lstlisting}[style=bash]
./prog < test1.in > test1.out
diff test1.out test1.expect
\end{lstlisting}

\subsection{Testing modules: a test client}

\begin{lstlisting}[style=c]
// test-stack.c: assertion-based tests for the stack module
#include <assert.h>
#include "stack.h"

int main(void) {
  struct stack s;
  stack_init(&s);
  assert(stack_is_empty(&s));       // empty case
  stack_push(&s, 1);
  assert(!stack_is_empty(&s));
  assert(stack_top(&s) == 1);       // one item
  stack_push(&s, 2);
  stack_push(&s, 3);
  assert(stack_pop(&s) == 3);       // LIFO order
  assert(stack_pop(&s) == 2);
  assert(stack_pop(&s) == 1);
  assert(stack_is_empty(&s));       // back to empty
}
\end{lstlisting}
An I/O-driven test harness is also common: \code{main} reads commands like
\code{push 5} / \code{pop} / \code{top} and prints results, so many tests can be written as
\code{.in}/\code{.expect} files without recompiling.

\subsection{Debugging strategies}
\begin{itemize}
  \item Read the error message: assertion failures give the file and line.
  \item Make the failing test as small as possible.
  \item Add tracing output (with a \code{DEBUG} flag) around the suspected code.
  \item Trace the memory model by hand: draw the stack frames and pointers.
  \item Check the usual suspects: uninitialized variables, off-by-one loop bounds, \code{=} vs
    \code{==}, missing \code{\&} in \code{scanf}, integer division.
\end{itemize}

\section{Summary}
\begin{itemize}
  \item \code{printf} writes to stdout; use a \code{DEBUG} flag for tracing output.
  \item \code{scanf} needs pointers and returns the number of items read (or \code{EOF}); always
    check it. \code{" \%c"} skips whitespace.
  \item Ctrl+D sends EOF in Seashell; \code{getchar} returns an \code{int} so \code{EOF} fits.
  \item Test black-box and white-box; cover edge cases; use \code{.in}/\code{.expect} files for
    I/O tests and \code{assert}-based test clients for modules; re-run all tests after changes.
\end{itemize}
